\section{Experimental Setup}\label{sec:setup}

\paragraph{Model.}
All experiments run on our from-scratch reproduction of Qwen3-0.6B-Base
\citep{yang2025qwen3}: a pre-norm, decoder-only transformer
\citep{vaswani2017attention} with 28 layers and hidden size 1{,}024.
Attention is grouped-query \citep{ainslie2023gqa} with 16 query heads and
8 key--value heads at an explicit head dimension of 128, set independently
in the configuration (it is not hidden size divided by head count, which
would give 64). Each block applies RMSNorm \citep{zhang2019root} with
$\epsilon = 10^{-6}$ before attention and before a SwiGLU feed-forward
network \citep{shazeer2020glu} of intermediate size 3{,}072; Qwen3
additionally normalizes queries and keys per head (QK-norm, an RMSNorm
over the 128-dimensional head axis) before rotary position embeddings
\citep{su2021roformer} with base $\theta = 10^{6}$ and positions up to
40{,}960. All projections are bias-free, and the input embedding is tied
to the output head over a 151{,}936-entry vocabulary. Recomputed from
this configuration, the model has exactly 596{,}049{,}920 parameters:
28 blocks of 15{,}730{,}944 each, a tied embedding matrix of
155{,}582{,}464, and a 1{,}024-parameter final norm
(Table~\ref{tab:repro}). A second reproduction, SmolLM2-135M
\citep{allal2025smollm2} with 134{,}515{,}008 parameters, shows that the
verification protocol is not specific to one model family
(Table~\ref{tab:repro}).

\paragraph{Tokenizer.}
We use the Qwen3 byte-level BPE tokenizer with vocabulary size
151{,}936. Every perplexity in this paper is computed with the model's
own tokenizer; cross-tokenizer perplexity comparisons are never made.
Cross-corpus effects are reported in bits per byte (BPB), which
normalizes cross-entropy by the UTF-8 byte count of the scored text and
is tokenizer-independent.

\paragraph{Data.}
The pre-training corpus is FineWeb-Edu \citep{penedo2024fineweb}
(\texttt{sample-10BT}), streamed and tokenized on the fly. The baseline
consumes 1{,}189{,}478{,}400 training tokens (18{,}150 steps at
65{,}536 tokens per step; Table~\ref{tab:hparams}), approximately
2 tokens per parameter -- deliberately far below the compute-optimal
ratio \citep{hoffmann2022training}, since the study optimizes for
controlled comparisons rather than absolute quality. A held-out
validation slice of 300{,}012 tokens is document-disjoint from training
via a seeded hash split and screened by 13-gram overlap decontamination
(0 of 451 validation documents dropped); the in-loop evaluator scores
50 windows of 4{,}096 tokens (204{,}800 tokens) on this slice. The
data-composition experiments add a dclm-edu slice \citep{li2024datacomp}
of 150{,}107{,}480 training tokens, pinned by content hash
(\texttt{dbad8ad7}) and decontaminated with the same 13-gram screen
against both evaluation corpora (0 documents dropped). Supervised
fine-tuning uses OpenR1-Math-220k pairs totalling 125{,}000{,}592
tokens. The GRPO prompt set (\texttt{grpo-math-prompts-v1}) contains
7{,}471 GSM8K and 3{,}490 MATH level 1--3 training prompts after
dropping 14 prompts that overlapped the evaluation sets at the 13-gram
level; 7 prompts overlapping the SFT set were flagged and kept. The math
evaluation sets themselves (\texttt{math-eval-v1}) are decontaminated
against all 35{,}924 OpenR1 SFT problems: 0 of 1{,}319 GSM8K test items
and 0 of 500 MATH-500 items are flagged.

\paragraph{Hardware.}
Everything runs on one NVIDIA GB10 (Grace Blackwell) workstation whose
CPU and GPU share a single $\approx$119 GB unified-memory pool with no
separate VRAM. Training uses bfloat16 with \texttt{torch.compile} at
sequence length 4{,}096; cross-entropy over the 151{,}936-entry
vocabulary is computed in fp32 chunks, since materializing the full fp32
logit matrix at this vocabulary size is infeasible in the shared pool.
A 50-step probe measured 7{,}167.4 tokens/s compiled versus 3{,}787.5
uncompiled (a 1.89$\times$ speedup), and the full baseline run sustained
7{,}444--7{,}482 tokens/s at a flat 52.4 GB peak memory, completing in
2{,}663.1 minutes. FLOP accounting follows the PaLM convention
\citep{chowdhery2022palm}. Because the GB10's peak throughput is an
estimate rather than a published figure, model FLOPs utilization is
reported as approximate only and never as an exact number.

\paragraph{Evaluation protocol.}
Three versioned suites produce every cross-run comparable number, and
each such number carries its suite stamp. (i) \textbf{text-lm-v2}:
perplexity and BPB on the wikitext-2-raw-v1 validation split and
codeparrot-clean-valid, both at pinned dataset revisions
(\texttt{b08601e}, \texttt{4db92d2}), scoring 204{,}600 tokens per
corpus per cell. The harness measures a per-corpus noise floor on the
baseline checkpoint: 1.304 PPL (3.73\% relative) on wikitext-2 but
280.6 PPL (68.2\% relative) on the code corpus (Table~\ref{tab:arms}).
Code-perplexity orderings are therefore always quoted with this caveat,
and comparative effects are instead reported as BPB improvements with
paired 3-seed 95\% confidence intervals, where seeds vary initialization
and data order; deltas below the floor or with an interval crossing zero
are labeled not significant. (ii) \textbf{math-acc-v1}: a pinned answer
extractor and verifier over GSM8K \citep{cobbe2021training} and MATH-500
levels 1--3 \citep{hendrycks2021measuring}, with pass@$k$ computed by
the unbiased estimator of \citet{chen2021evaluating} and Wilson 95\%
intervals for pass@1, under a fixed sampling band (8 samples per item,
temperature 0.8, 256 new tokens, fixed band seed) on 100 GSM8K and
50 MATH-500 items. (iii) \textbf{ecl-ladder-v1}: passkey retrieval
\citep{mohtashami2023landmark} across context rungs 512--8{,}192 at
depths 0.1--0.9, with 40 probes per rung per cell and Wilson intervals,
an effective-context-length measurement in the spirit of RULER
\citep{hsieh2024ruler}. Trainer-internal perplexities are labeled
throughout as in-loop validation perplexity, single run, not
suite-stamped, and are advisory only. Comparative arms are
single-variable and iso-FLOP by construction: two arms are compared only
if exactly one variable differs and their training FLOPs match within
5\%, with matched token and step budgets stated for every comparison.

\paragraph{Hyperparameters.}
Table~\ref{tab:hparams} consolidates the per-stage hyperparameters --
optimizer, learning-rate schedule, warmup, weight decay, batch shape,
token and step budgets, and seeds -- with each value flagged as reported
in its source paper, inferred from a published band, or chosen by us
where no published value exists at this scale. AdamW
\citep{loshchilov2017decoupled} is the optimizer at every stage except
the treatment arm of the optimizer A/B itself.

% ---------------------------------------------------------------------
% tab_hparams.tex -- per-stage hyperparameters for every experiment
% stage in the paper. Every value is transcribed from the run's
% c5_evidence.json recipe, the trainer's argparse defaults, or the
% args/start line of the on-disk training log.
% Provenance flags: R = reported in the cited source paper;
% I = inferred (source gives a band, we picked a point);
% C = chosen by us (no published value at this scale/budget).
% ASCII only; requires only booktabs.
% ---------------------------------------------------------------------
\providecommand{\hpR}{\textsuperscript{R}}
\providecommand{\hpI}{\textsuperscript{I}}
\providecommand{\hpC}{\textsuperscript{C}}

\begin{table}[t]
\centering
\caption{Per-stage training hyperparameters. Superscripts flag
provenance: \hpR{} reported in the source paper (Qwen3 technical
report LR shape; VibeThinker SFT stage-1; GRPO/DAPO/Dr.GRPO),
\hpI{} inferred from a band in the source, \hpC{} chosen by us (no
published value at this scale). Shared unless stated: AdamW betas
$(0.9,0.95)$, eps $10^{-8}$, grad-clip 1.0, bf16, seq.\ len.\ 4096,
chunked fp32 cross-entropy, one GB10 GPU; all part-(a) stages use
65{,}536 tok/step (micro-batch 4 $\times$ grad-accum 4)\hpC. Each run
directory records its full recipe in \texttt{c5\_evidence.json} before
launch (contract C5); values are cross-checked against the
\texttt{args} lines of the training logs. The optimizer A/B also ran a
matched-config AdamW LR sweep $\{1.7,2.4,3.5,4.8\}\times10^{-3}$ (n=1
off-baseline). GRPO Phase 2 is n=1 by design (verdict capped at
directional).}
\label{tab:hparams}
\footnotesize
\setlength{\tabcolsep}{3.5pt}
\renewcommand{\arraystretch}{1.05}

% ---- (a) Pre-training and mid-training stages ------------------------
\begin{tabular}{@{}lp{2.0cm}p{2.2cm}p{2.6cm}p{2.3cm}p{2.2cm}@{}}
\toprule
 & \textbf{Faithful 2TPP} & \textbf{Optimizer A/B} & \textbf{De-confound P1/P2} & \textbf{Data cells} & \textbf{Mid-train anneal} \\
 & (06-08) & (06-16) & (06-18 / 06-21) & (06-24 / 06-26) & (06-30) \\
\midrule
Init            & random & random & random & random & 2TPP ckpt \\
Data            & FineWeb-Edu\hpC & FineWeb-Edu\hpC & FineWeb-Edu\hpC & dclm-edu; 50/50 mix\hpC & dclm-edu, FineWeb\hpC \\
Arms            & --- & AdamW vs NorMuon-2D & +WSD, +z-loss, +arch; VR, LNS, HG & dclm; mix (ctrl reused) & 50/50 mix vs FineWeb-only \\
Optimizer       & AdamW\hpR & AdamW vs NorMuon\hpC & AdamW\hpR & AdamW\hpR & AdamW\hpR \\
LR peak $\to$ end & 2.4e-3\hpC $\to$ 3.2e-4\hpR & 2.4e-3\hpC; NM 0.011\hpC{} ($\to$ 0.1x)\hpC & 1.7e-3\hpR $\to$ 3.2e-4\hpR & 1.7e-3\hpR $\to$ 3.2e-4\hpR & 2.5e-4\hpC $\to$ 2.5e-5\hpC \\
Schedule        & cosine\hpR & cosine\hpR & cosine\hpR{} (WSD arm: WSD) & cosine\hpR & 1-sqrt decay\hpR{} (shape) \\
Warmup (steps)  & 900\hpC & 50\hpC & 100\hpC & 100\hpC & 46\hpC \\
Weight decay    & 0.01\hpR & 0.1 (2D)\hpC & 0.01\hpR & 0.01\hpR & 0.01\hpR \\
Steps (tok)/cell & 18{,}150 (1.19B, 2~TPP\hpC) & 640 (41.9M\hpC) & 2{,}000 (131M\hpC) & 2{,}000 (131M\hpC) & 2{,}300 (151M\hpC) \\
Cells (new)     & 1 & 6 & 12 / 9 & 3 / 3 & 6 \\
Seeds           & 0 & 0--2 & 0--2 & 0--2 & 0--2 \\
\bottomrule
\end{tabular}

\vspace{5pt}

% ---- (b) Post-training stages ----------------------------------------
\begin{tabular}{@{}lp{4.6cm}p{5.6cm}@{}}
\toprule
 & \textbf{SFT 3-seed} (06-27) & \textbf{GRPO Phase 2} (07-02) \\
\midrule
Init             & faithful 2TPP base (strict load) & SFT seed-0 ckpt (policy + frozen KL ref) \\
Data             & OpenR1-Math-220k pairs\hpC & grpo-math-prompts-v1: GSM8K + MATH L1--3, 2:1 mix\hpC \\
Arms             & masked SFT vs no-mask iso-FLOP control & Dr.GRPO vs random-reward vs RFT \\
Optimizer / LR   & AdamW, 5e-5\hpR $\to$ 8e-8\hpR{} cosine & AdamW, 1e-6 const.\hpI{} (RFT: 1e-5\hpI) \\
Warmup / wd      & 5\% (11 steps)\hpR{} / 0.01\hpC & none / 0.0\hpC \\
Batch            & 524{,}288 tok/step (mb 4 x ga 32 = 128 seqs\hpR) & 16 prompts x G=8 rollouts / step\hpC \\
Steps (budget)   & 238 (125M tok\hpC) & 300 (300 x 16 x 8 rollouts, cap 256 new tok\hpC) \\
Sampling         & --- & T = 0.8\hpC, no KV cache (exact re-forward) \\
Clip / KL        & --- & clip eps 0.2\hpR; k3-KL beta 3e-3\hpI \\
Reward           & --- & math-acc-v1: 1.0 / 0.1 / 0.0\hpC; random arm Bernoulli(0.5) \\
Seeds            & 0--2 & 0 (n=1, directional) \\
\bottomrule
\end{tabular}
\end{table}

